% =============================================================================
% CAPÍTULO 1: INTRODUÇÃO
% =============================================================================

\chapter{Introdução}

\section{Contextualização}

As máquinas rotativas constituem elementos fundamentais em diversos setores industriais, desde a geração de energia até processos de manufatura. Sua operação confiável é essencial para a eficiência e segurança dos processos industriais, sendo que falhas não detectadas podem resultar em paradas não planejadas, perdas de produção e, em casos extremos, acidentes com consequências graves.

O monitoramento e diagnóstico de falhas em máquinas rotativas tem evoluído significativamente nas últimas décadas, impulsionado pelo desenvolvimento de técnicas avançadas de análise de sinais e pela crescente demanda por sistemas de manutenção preditiva. Neste contexto, a análise de vibrações emerge como uma das principais ferramentas para o diagnóstico de falhas, oferecendo informações valiosas sobre o estado operacional dos equipamentos.

\section{Justificativa}

A análise de espectros de ordem superior (Higher-Order Spectra - HOS), incluindo bi-espectro e tri-espectro, oferece uma abordagem promissora para identificar assinaturas não lineares associadas a diferentes tipos de falhas. No entanto, a aplicação prática dessas técnicas muitas vezes enfrenta desafios devido à complexidade dos fenômenos envolvidos e à necessidade de dados experimentais extensivos.

Este trabalho propõe o desenvolvimento de uma metodologia baseada em simulações numéricas para o diagnóstico de falhas em máquinas rotativas, utilizando a análise de HOS. A abordagem por simulação oferece vantagens significativas, incluindo controle preciso dos parâmetros do sistema, capacidade de gerar grandes volumes de dados e possibilidade de estudar cenários que seriam difíceis ou perigosos de reproduzir experimentalmente.

\section{Problema de Pesquisa}

O problema central desta pesquisa consiste em desenvolver uma metodologia eficaz para o diagnóstico de falhas em máquinas rotativas utilizando simulações numéricas e análise de espectros de ordem superior, com o objetivo de identificar padrões característicos que permitam distinguir entre condições normais de operação e diferentes tipos de falhas.

\section{Objetivos}

\subsection{Objetivo Geral}

Desenvolver uma metodologia baseada em simulações numéricas para o diagnóstico de falhas em máquinas rotativas utilizando análise de espectros de ordem superior.

\subsection{Objetivos Específicos}

\begin{enumerate}
    \item Desenvolver modelos de elementos finitos utilizando a ferramenta ROSS para simular o comportamento dinâmico de eixos rotativos sob condições normais e com falhas específicas, como trincas e desalinhamentos.
    
    \item Aplicar técnicas de análise de espectros de ordem superior (bi-espectro e tri-espectro) aos dados simulados para identificar padrões característicos que diferenciem as condições normais das falhas modeladas.
    
    \item Validar os resultados obtidos através de comparação com dados disponíveis na literatura científica.
    
    \item Avaliar a eficácia da metodologia proposta para diferentes tipos de falhas e condições operacionais.
\end{enumerate}

\section{Organização do Trabalho}

Esta dissertação está organizada em cinco capítulos, além dos elementos pré e pós-textuais:

\begin{itemize}
    \item \textbf{Capítulo 1 - Introdução:} Apresenta a contextualização, justificativa, problema de pesquisa, objetivos e organização do trabalho.
    
    \item \textbf{Capítulo 2 - Revisão de Literatura:} Aborda os fundamentos teóricos relacionados à dinâmica de rotores, análise de vibrações, espectros de ordem superior e diagnóstico de falhas.
    
    \item \textbf{Capítulo 3 - Metodologia:} Descreve detalhadamente a metodologia proposta, incluindo a modelagem numérica, simulações e análise de dados.
    
    \item \textbf{Capítulo 4 - Resultados e Discussão:} Apresenta e discute os principais resultados obtidos, incluindo a análise dos espectros de ordem superior e a identificação de padrões característicos.
    
    \item \textbf{Capítulo 5 - Conclusões:} Sintetiza os principais achados, contribuições e perspectivas futuras do trabalho.
\end{itemize}
